% Options for packages loaded elsewhere
\PassOptionsToPackage{unicode}{hyperref}
\PassOptionsToPackage{hyphens}{url}
\documentclass[
  12pt,
]{ctexart}
\usepackage{xcolor}
\usepackage{amsmath,amssymb}
\setcounter{secnumdepth}{-\maxdimen} % remove section numbering
\usepackage{iftex}
\ifPDFTeX
  \usepackage[T1]{fontenc}
  \usepackage[utf8]{inputenc}
  \usepackage{textcomp} % provide euro and other symbols
\else % if luatex or xetex
  \usepackage{unicode-math} % this also loads fontspec
  \defaultfontfeatures{Scale=MatchLowercase}
  \defaultfontfeatures[\rmfamily]{Ligatures=TeX,Scale=1}
\fi
\usepackage{lmodern}
\ifPDFTeX\else
  % xetex/luatex font selection
\fi
% Use upquote if available, for straight quotes in verbatim environments
\IfFileExists{upquote.sty}{\usepackage{upquote}}{}
\IfFileExists{microtype.sty}{% use microtype if available
  \usepackage[]{microtype}
  \UseMicrotypeSet[protrusion]{basicmath} % disable protrusion for tt fonts
}{}
\makeatletter
\@ifundefined{KOMAClassName}{% if non-KOMA class
  \IfFileExists{parskip.sty}{%
    \usepackage{parskip}
  }{% else
    \setlength{\parindent}{0pt}
    \setlength{\parskip}{6pt plus 2pt minus 1pt}}
}{% if KOMA class
  \KOMAoptions{parskip=half}}
\makeatother
\usepackage{longtable,booktabs,array}
\usepackage{calc} % for calculating minipage widths
% Correct order of tables after \paragraph or \subparagraph
\usepackage{etoolbox}
\makeatletter
\patchcmd\longtable{\par}{\if@noskipsec\mbox{}\fi\par}{}{}
\makeatother
% Allow footnotes in longtable head/foot
\IfFileExists{footnotehyper.sty}{\usepackage{footnotehyper}}{\usepackage{footnote}}
\makesavenoteenv{longtable}
\setlength{\emergencystretch}{3em} % prevent overfull lines
\providecommand{\tightlist}{%
  \setlength{\itemsep}{0pt}\setlength{\parskip}{0pt}}
\usepackage{bookmark}
\IfFileExists{xurl.sty}{\usepackage{xurl}}{} % add URL line breaks if available
\urlstyle{same}
\hypersetup{
  hidelinks,
  pdfcreator={LaTeX via pandoc}}

\author{SCX}
\date{2026}

\begin{document}

\section{SCX 整理：总体框架、模块划分、文件夹结构、论文分布与 LLM
定位}\label{scx-ux6574ux7406ux603bux4f53ux6846ux67b6ux6a21ux5757ux5212ux5206ux6587ux4ef6ux5939ux7ed3ux6784ux8bbaux6587ux5206ux5e03ux4e0e-llm-ux5b9aux4f4d}

\begin{quote}
基于 \texttt{与GPT的讨论2026-06-26-07-49.md} (9011 行) 的系统整理
整理日期：2026-06-26
\end{quote}

\begin{center}\rule{0.5\linewidth}{0.5pt}\end{center}

\subsection{一、SCX
总体框架定义}\label{ux4e00scx-ux603bux4f53ux6846ux67b6ux5b9aux4e49}

\subsubsection{1.1 两层架构}\label{ux4e24ux5c42ux67b6ux6784}

GPT 明确建议 SCX 走 \textbf{``论文确权 + 最小开源 + 核心工程保留 +
可商业授权''} 的路线，即两层架构：

\paragraph{第一层：SCX-Research（开源层）}\label{ux7b2cux4e00ux5c42scx-researchux5f00ux6e90ux5c42}

\begin{longtable}[]{@{}ll@{}}
\toprule\noalign{}
维度 & 说明 \\
\midrule\noalign{}
\endhead
\bottomrule\noalign{}
\endlastfoot
\textbf{定位} & 学术层、开源核心 \\
\textbf{目标用户} & 学术界、开源社区、审稿人、潜在合作者 \\
\textbf{许可证} & 偏宽松（便于学术传播） \\
\textbf{开放程度} & 完全开源 \\
\textbf{目的} & 让别人能复现论文、建立学术影响力、抢概念命名权 \\
\end{longtable}

内容清单： - SCX 数学定义与核心公式 - 命题与证明 - 算法伪代码 -
Synthetic experiments（合成实验） - Toy data valuation demo -
State-conditioned expert reliability estimator（基础版） - Simple
redundancy compression（基础版） - Paper figures reproduction - 一个小型
MLIP / 通用 ML demo

\paragraph{第二层：SCX-Pro /
SCX-Compiler（闭源/商业层）}\label{ux7b2cux4e8cux5c42scx-pro-scx-compilerux95edux6e90ux5546ux4e1aux5c42}

\begin{longtable}[]{@{}ll@{}}
\toprule\noalign{}
维度 & 说明 \\
\midrule\noalign{}
\endhead
\bottomrule\noalign{}
\endlastfoot
\textbf{定位} & 商业执行层、工程实现 \\
\textbf{目标用户} & 企业客户、产业联盟成员 \\
\textbf{开放程度} & 闭源，通过授权/年费/项目费交付 \\
\textbf{目的} & 形成商业壁垒，服务付费客户 \\
\end{longtable}

内容清单： - 大规模数据管线（large-scale data pipeline） - MLIP expert
compiler（多势函数编译/蒸馏） - LLM data curation
module（大模型数据治理） - 企业可视化 dashboard -
自动报告生成（automated report generation） - 专家库管理（expert
registry） - 企业私有部署 / API 服务 - 私有 benchmark -
工业级自动阈值选择/状态划分策略 - 专家冲突仲裁完整工程细节

\paragraph{第三层（远期）：SCX
Consortium}\label{ux7b2cux4e09ux5c42ux8fdcux671fscx-consortium}

产业联盟层： - 多方公司交会员费参与共建 - 核心标准由 SCX 维护方控制 -
非独占授权，不卖断 - 形成多方认可的中立标准

\subsubsection{1.2 核心原则}\label{ux6838ux5fc3ux539fux5219}

\begin{quote}
\textbf{理论开放，基础代码开放，工程系统保留。}
\end{quote}

\begin{itemize}
\tightlist
\item
  论文保护学术优先权
\item
  开源保护生态扩散
\item
  专利保护商业排他性
\end{itemize}

\begin{quote}
\textbf{开源到''足以复现论文''，不要开源到''足以替代商业产品''。}
\end{quote}

\subsubsection{1.3 IP 保护建议}\label{ip-ux4fddux62a4ux5efaux8bae}

顺序链： 1. 先做 IP 评估/学校成果披露 2. 有专利价值先提交申请 3.
再发论文/arXiv 4. 再开源最小代码 5. 最后做商业版

\begin{center}\rule{0.5\linewidth}{0.5pt}\end{center}

\subsection{二、模块划分}\label{ux4e8cux6a21ux5757ux5212ux5206}

\subsubsection{2.1
核心模块（SCX-Core）}\label{ux6838ux5fc3ux6a21ux5757scx-core}

这些是 SCX 的数学/算法基础，属于开源层。

\begin{longtable}[]{@{}
  >{\raggedright\arraybackslash}p{(\linewidth - 6\tabcolsep) * \real{0.2143}}
  >{\raggedright\arraybackslash}p{(\linewidth - 6\tabcolsep) * \real{0.3571}}
  >{\raggedright\arraybackslash}p{(\linewidth - 6\tabcolsep) * \real{0.2143}}
  >{\raggedright\arraybackslash}p{(\linewidth - 6\tabcolsep) * \real{0.2143}}@{}}
\toprule\noalign{}
\begin{minipage}[b]{\linewidth}\raggedright
模块
\end{minipage} & \begin{minipage}[b]{\linewidth}\raggedright
功能定义
\end{minipage} & \begin{minipage}[b]{\linewidth}\raggedright
输入
\end{minipage} & \begin{minipage}[b]{\linewidth}\raggedright
输出
\end{minipage} \\
\midrule\noalign{}
\endhead
\bottomrule\noalign{}
\endlastfoot
\textbf{State Discovery（状态发现）} & 提取 embedding，聚类形成状态空间
& 数据集 D = \{(x\_i, y\_i)\}，表示函数 phi(x) & 状态标签
s\_i，状态边界 \\
\textbf{Expert Reliability（专家可靠性估计）} &
估计每个专家在每个状态上的可靠性 & 多专家输出 f\_m(x)，gold labels y*
或代理标签 & SCX\_m(s) = P(loss \textless{} tau \\
\textbf{Data Value Function（数据价值函数）} & 计算样本/状态的数据价值 &
状态统计量：mean residual, density, similarity, consistency & V(s)
价值分数 \\
\textbf{Redundancy Compression（冗余压缩）} & 识别低误差高密度状态并压缩
& 状态统计量、embedding & 代表样本子集 + 权重 \\
\textbf{Noise Risk Estimation（噪声风险评估）} &
估计噪声风险，区分可学习困难状态和噪声 & residual, density, consistency,
anchor support & N(x) 噪声风险分数 \\
\textbf{Expert Routing（专家路由）} & 按状态选择最优专家 & 各专家可靠性
R\_m(s)，成本 C\_m & m*(x) = argmin{[}R\_m(s) + lambda*C\_m{]} \\
\textbf{Action Policy（数据动作策略）} & 输出最终数据操作决策 &
上述所有模块输出 & keep/drop/downweight/relabel/route \\
\end{longtable}

\subsubsection{2.2 领域适配模块（Domain
Adapters）}\label{ux9886ux57dfux9002ux914dux6a21ux5757domain-adapters}

这些是 SCX-Core 在各行业的具体实现。

\paragraph{A. SCX-MLIP（材料势函数 ---
主战场/根据地）}\label{a.-scx-mlipux6750ux6599ux52bfux51fdux6570-ux4e3bux6218ux573aux6839ux636eux5730}

\begin{longtable}[]{@{}
  >{\raggedright\arraybackslash}p{(\linewidth - 2\tabcolsep) * \real{0.5000}}
  >{\raggedright\arraybackslash}p{(\linewidth - 2\tabcolsep) * \real{0.5000}}@{}}
\toprule\noalign{}
\begin{minipage}[b]{\linewidth}\raggedright
维度
\end{minipage} & \begin{minipage}[b]{\linewidth}\raggedright
说明
\end{minipage} \\
\midrule\noalign{}
\endhead
\bottomrule\noalign{}
\endlastfoot
\textbf{状态表示} & ACE descriptor / atomic local environment \\
\textbf{参考标签} & DFT (energy, force, stress) \\
\textbf{专家类型} & ACE, NEP, MACE, DeepMD, GAP, NequIP 等势函数 \\
\textbf{核心功能} & descriptor residual map, gauge alignment,
energy-zero alignment, expert merging, DFT budget-aware active
learning \\
\textbf{数据动作} & 补 DFT / 合并专家 / 蒸馏学生 / 压缩冗余 AIMD 帧 \\
\textbf{商业价值} & 势函数数据审计、多专家编译蒸馏 \\
\end{longtable}

\paragraph{B. SCX-Sim / SCX-FEM（工程仿真 ---
商业主入口）}\label{b.-scx-sim-scx-femux5de5ux7a0bux4effux771f-ux5546ux4e1aux4e3bux5165ux53e3}

\begin{longtable}[]{@{}
  >{\raggedright\arraybackslash}p{(\linewidth - 2\tabcolsep) * \real{0.5000}}
  >{\raggedright\arraybackslash}p{(\linewidth - 2\tabcolsep) * \real{0.5000}}@{}}
\toprule\noalign{}
\begin{minipage}[b]{\linewidth}\raggedright
维度
\end{minipage} & \begin{minipage}[b]{\linewidth}\raggedright
说明
\end{minipage} \\
\midrule\noalign{}
\endhead
\bottomrule\noalign{}
\endlastfoot
\textbf{状态表示} & 几何/网格/响应 embedding, 输入参数 \\
\textbf{参考标签} & 细网格/高保真 FEM / 实验数据 \\
\textbf{专家类型} & 粗网格 FEM, 细网格 FEM, 高阶求解器, ROM, neural
surrogate \\
\textbf{核心功能} &
识别高价值工况、压缩冗余仿真、识别数值异常、多保真专家路由 \\
\textbf{数据动作} & 高保真重算 / 保留 / 压缩 / 异常标记 \\
\textbf{商业价值} & 仿真数据审计报告、surrogate
训练集选择、多保真计算调度 \\
\end{longtable}

\paragraph{C. SCX-Health（医学/可穿戴健康 ---
开源影响力）}\label{c.-scx-healthux533bux5b66ux53efux7a7fux6234ux5065ux5eb7-ux5f00ux6e90ux5f71ux54cdux529b}

\begin{longtable}[]{@{}
  >{\raggedright\arraybackslash}p{(\linewidth - 2\tabcolsep) * \real{0.5000}}
  >{\raggedright\arraybackslash}p{(\linewidth - 2\tabcolsep) * \real{0.5000}}@{}}
\toprule\noalign{}
\begin{minipage}[b]{\linewidth}\raggedright
维度
\end{minipage} & \begin{minipage}[b]{\linewidth}\raggedright
说明
\end{minipage} \\
\midrule\noalign{}
\endhead
\bottomrule\noalign{}
\endlastfoot
\textbf{状态表示} & 图像 embedding (DINO/CLIP/ViT)，临床特征 \\
\textbf{参考标签} & 医生标签 / 多医生一致性 / 金标准 \\
\textbf{专家类型} & 不同模型、不同医生、不同影像设备 \\
\textbf{核心功能} &
冗余正常病例压缩、噪声标签识别、长尾高价值病例、专家复核优先排序 \\
\textbf{数据动作} & 保留 / 压缩 / 降权 / 专家复核 / 补标注 \\
\textbf{商业价值} & 第三方数据价值审计报告、可穿戴健康数据治理 \\
\end{longtable}

\paragraph{D. SCX-Semi（半导体仿真/制造 ---
远期战略）}\label{d.-scx-semiux534aux5bfcux4f53ux4effux771fux5236ux9020-ux8fdcux671fux6218ux7565}

\begin{longtable}[]{@{}
  >{\raggedright\arraybackslash}p{(\linewidth - 2\tabcolsep) * \real{0.5000}}
  >{\raggedright\arraybackslash}p{(\linewidth - 2\tabcolsep) * \real{0.5000}}@{}}
\toprule\noalign{}
\begin{minipage}[b]{\linewidth}\raggedright
维度
\end{minipage} & \begin{minipage}[b]{\linewidth}\raggedright
说明
\end{minipage} \\
\midrule\noalign{}
\endhead
\bottomrule\noalign{}
\endlastfoot
\textbf{子模块} & SCX-CMP, SCX-TCAD, SCX-VM (Virtual Metrology),
SCX-OPC, SCX-Yield \\
\textbf{状态表示} & layout density map, pattern clip embedding, device
geometry, process parameters \\
\textbf{专家类型} & Rule-based OPC, Model-based OPC, ILT, ML-OPC,
粗网格/细网格 TCAD, 简化/工业 CMP 模型 \\
\textbf{核心功能} & hotspot 可信度校准、多保真调度、候选 mask
排序、layout 修改优先排序、VM 预测可信度评估 \\
\textbf{数据动作} & 普通 OPC / ILT / 高保真仿真 / 量测 / 重算 \\
\textbf{商业价值} & EDA/TCAD 上层质量判断层 \\
\end{longtable}

\paragraph{E.
SCX-Vision（工业/通用视觉）}\label{e.-scx-visionux5de5ux4e1aux901aux7528ux89c6ux89c9}

\begin{longtable}[]{@{}
  >{\raggedright\arraybackslash}p{(\linewidth - 2\tabcolsep) * \real{0.5000}}
  >{\raggedright\arraybackslash}p{(\linewidth - 2\tabcolsep) * \real{0.5000}}@{}}
\toprule\noalign{}
\begin{minipage}[b]{\linewidth}\raggedright
维度
\end{minipage} & \begin{minipage}[b]{\linewidth}\raggedright
说明
\end{minipage} \\
\midrule\noalign{}
\endhead
\bottomrule\noalign{}
\endlastfoot
\textbf{状态表示} & DINO/CLIP/ViT embedding, defect morphology \\
\textbf{参考标签} & 人工标注 / 质检标签 \\
\textbf{核心功能} &
冗余正常样本压缩、高价值长尾缺陷、标注噪声识别、专家复核优先级 \\
\textbf{商业价值} & 工业质检数据治理、视觉标注成本降低 \\
\end{longtable}

\paragraph{F. SCX-Music（文化创意 ---
探索性/演示用）}\label{f.-scx-musicux6587ux5316ux521bux610f-ux63a2ux7d22ux6027ux6f14ux793aux7528}

\begin{longtable}[]{@{}ll@{}}
\toprule\noalign{}
维度 & 说明 \\
\midrule\noalign{}
\endhead
\bottomrule\noalign{}
\endlastfoot
\textbf{定位} & 有趣 demo，展示 SCX 泛化能力，非 Nature 主线 \\
\textbf{状态表示} & 音频特征 + 歌词特征 + 结构特征 \\
\textbf{参考标签} & 流媒体表现 (stream/save/share/skip) \\
\textbf{核心功能} & 歌曲状态分类 + 热歌条件规律 + demo 质量诊断 \\
\textbf{商业价值} & demo 评分器、A\&R 筛歌系统、热歌状态地图 \\
\textbf{注意} & 音乐标签极脏，受市场/平台/歌手/运气等多因素干扰 \\
\end{longtable}

\paragraph{G. SCX-LLM / SCX-Judge（大模型数据治理 ---
远期外延）}\label{g.-scx-llm-scx-judgeux5927ux6a21ux578bux6570ux636eux6cbbux7406-ux8fdcux671fux5916ux5ef6}

\begin{longtable}[]{@{}
  >{\raggedright\arraybackslash}p{(\linewidth - 2\tabcolsep) * \real{0.5000}}
  >{\raggedright\arraybackslash}p{(\linewidth - 2\tabcolsep) * \real{0.5000}}@{}}
\toprule\noalign{}
\begin{minipage}[b]{\linewidth}\raggedright
维度
\end{minipage} & \begin{minipage}[b]{\linewidth}\raggedright
说明
\end{minipage} \\
\midrule\noalign{}
\endhead
\bottomrule\noalign{}
\endlastfoot
\textbf{定位} & 评价器编排框架，非主战场 \\
\textbf{状态表示} & task embedding / 领域 embedding / judge
disagreement \\
\textbf{专家类型} & LLM judge, reward model, verifier, influence score,
loss dynamics, DSIR \\
\textbf{核心功能} & 多评价器状态条件加权、judge
可靠性校准、高分冗余/低分高价值判断 \\
\textbf{数据动作} & keep/drop/downweight/relabel/route to stronger
judge/verify with tool \\
\end{longtable}

\subsubsection{2.3
模块间依赖关系}\label{ux6a21ux5757ux95f4ux4f9dux8d56ux5173ux7cfb}

\begin{verbatim}
SCX-Core (抽象层)
    |
    ├── phi_domain (领域状态编码器接口)
    │       ├── SCX-MLIP Encoder (ACE descriptor)
    │       ├── SCX-FEM Encoder (几何/网格/模态)
    │       ├── SCX-Vision Encoder (DINO/CLIP/ViT)
    │       ├── SCX-Semi Encoder (layout/process)
    │       ├── SCX-Music Encoder (audio/lyrics)
    │       └── SCX-LLM Encoder (task embedding)
    │
    ├── State Statistics (通用)
    │       ├── Residual estimation
    │       ├── Density estimation
    │       ├── Similarity / Redundancy
    │       └── Consistency / Expert reliability
    │
    ├── Data Action Policy (通用)
    │       ├── Compress / keep / drop / downweight
    │       ├── Expert routing
    │       └── Training budget allocation
    │
    └── Industry Applications
            ├── SCX-MLIP (根据地)
            ├── SCX-Sim/FEM (商业入口)
            ├── SCX-Health (开源影响力)
            ├── SCX-Semi (远期战略)
            ├── SCX-Vision (外延)
            ├── SCX-Music (demo)
            └── SCX-LLM (外延讨论)
\end{verbatim}

\begin{center}\rule{0.5\linewidth}{0.5pt}\end{center}

\subsection{三、文件夹结构设计}\label{ux4e09ux6587ux4ef6ux5939ux7ed3ux6784ux8bbeux8ba1}

\subsubsection{3.1
现有结构诊断}\label{ux73b0ux6709ux7ed3ux6784ux8bcaux65ad}

当前 \texttt{G:/Xiaogan\_Supercomputing\_data/SCX/} 结构：

\begin{verbatim}
SCX/
  src/scx/           # 核心代码
    core/            # 核心模块
    state/           # 状态发现
    expert/          # 专家模块
    valuation/       # 价值评估
    action/          # 动作策略
    utils/           # 工具函数
  theory/            # 数学定义
    definitions/
    propositions/
  paper/             # 论文
    paper_sources/
    paper_text/
    papers/
  experiments/       # 实验
    ml_benchmarks/
    mlip_case/
    synthetic/
  CodexKnowledge/    # 知识管理
\end{verbatim}

存在问题： - 没有显式区分 \texttt{scx-research/}（开源）和
\texttt{scx-pro/}（闭源/商业） - 领域适配器（MLIP, FEM, Health,
Vision）未在目录中体现 - 论文分布没有与模块对应

\subsubsection{3.2
建议的新文件夹布局}\label{ux5efaux8baeux7684ux65b0ux6587ux4ef6ux5939ux5e03ux5c40}

\begin{verbatim}
SCX/
├── README.md
├── pyproject.toml
│
├── scx-research/                  # 开源核心 (代码开源到"足以复现论文")
│   ├── core/                      # SCX 数学框架核心
│   │   ├── state_discovery.py     # 状态发现
│   │   ├── expert_reliability.py  # 专家可靠性估计
│   │   ├── data_value.py          # 数据价值函数
│   │   ├── redundancy_compress.py # 冗余压缩
│   │   ├── noise_risk.py          # 噪声风险评估
│   │   ├── expert_routing.py      # 专家路由
│   │   └── action_policy.py       # 数据动作策略
│   │
│   ├── adapters/                  # 领域适配器 (State Encoders)
│   │   ├── mlip/                  # SCX-MLIP
│   │   │   ├── ace_encoder.py     # ACE descriptor based state encoder
│   │   │   ├── gauge_align.py     # 能量规范对齐
│   │   │   └── residual_map.py    # 残差地图
│   │   ├── fem/                   # SCX-FEM/SCX-Sim
│   │   │   ├── fem_encoder.py     # FEM state encoder
│   │   │   ├── multi_fidelity.py  # 多保真路由
│   │   │   └── surrogate_select.py# surrogate 训练集选择
│   │   ├── vision/                # SCX-Vision
│   │   │   ├── vision_encoder.py  # DINO/CLIP based encoder
│   │   │   └── label_noise.py     # 标签噪声检测
│   │   ├── health/                # SCX-Health (单独也可作为独立仓库)
│   │   │   ├── datasets/          # 下载脚本、处理流程（不含原始受限数据）
│   │   │   ├── medmnist/          # MedMNIST 适配
│   │   │   ├── ham10000/          # HAM10000 适配
│   │   │   └── chexpert_lite/     # CheXpert 适配
│   │   ├── semi/                  # SCX-Semi (toy 版)
│   │   │   ├── cmp_toy/           # CMP toy model
│   │   │   ├── opc_toy/           # OPC toy lithography
│   │   │   └── tcad_like/         # TCAD-like PDE
│   │   ├── music/                 # SCX-Music (探索性)
│   │   └── llm/                   # SCX-LLM (讨论/外延)
│   │
│   ├── examples/                  # 使用示例
│   │   ├── synthetic/             # 合成数据实验
│   │   ├── small_classification/  # 小规模分类
│   │   ├── mlip_toy/              # MLIP toy demo
│   │   ├── medical_demo/          # 医学数据 demo
│   │   └── fem_demo/              # FEM demo
│   │
│   └── tests/                     # 单元测试
│
├── scx-pro/                       # 闭源/商业版 (模板/占位，不包含真实代码)
│   ├── README.md                  # 说明：商业版功能列表
│   ├── modules/                   # 模块说明（仅文档）
│   │   ├── large_scale_pipeline.md
│   │   ├── enterprise_dashboard.md
│   │   ├── expert_compiler.md
│   │   ├── llm_data_curation.md
│   │   └── auto_report.md
│   └── api/                       # API 接口定义（公开的接口规范）
│
├── experiments/                   # 实验记录/结果
│   ├── mlip_case/                 # MLIP 案例实验
│   ├── synthetic/                 # 合成实验
│   ├── ml_benchmarks/             # 通用 ML benchmark
│   ├── medical/                   # 医学实验记录
│   ├── fem/                       # FEM 实验记录
│   └── results/                   # 实验结果汇总
│       ├── figures/               # 论文用图
│       └── tables/                # 论文用表
│
├── paper/                         # 论文
│   ├── paper_0_egp/               # Paper 1: EGP — 实际位置在 egp/ (G:/Xiaogan/egp/)
│   │   └── README.md              # ACE gauge-normalized expert merging
│   ├── paper_1_theory/            # Paper 2: SCX-Theory (数学框架)
│   │   ├── draft/
│   │   ├── figures/
│   │   └── submission/
│   ├── paper_2_mlip/              # Paper 3: SCX-MLIP (理论→MLIP应用)
│   │   ├── draft/
│   │   ├── figures/
│   │   └── submission/
│   ├── paper_3_sim/               # Paper 4: SCX-Sim (多保真仿真)
│   │   ├── draft/
│   │   ├── figures/
│   │   └── submission/
│   ├── paper_4_health/            # Paper 5: SCX-Health (医学数据)
│   │   ├── draft/
│   │   ├── figures/
│   │   └── submission/
│   └── paper_common/              # 论文通用资源
│       ├── bibliography.bib
│       ├── templates/
│       └── style_files/
│
├── theory/                        # 数学理论 (独立于代码)
│   ├── definitions/               # 数学定义
│   └── propositions/              # 命题与证明
│
├── CodexKnowledge/                # 知识管理 (已有)
│   ├── 与GPT的讨论*.md
│   ├── 整理_框架与模块与论文分布.md (本文件)
│   ├── SCX_核心定义.md
│   └── ...
│
└── images/                        # 通用图片/架构图
\end{verbatim}

\subsubsection{3.3
开源/闭源分界清单}\label{ux5f00ux6e90ux95edux6e90ux5206ux754cux6e05ux5355}

\begin{longtable}[]{@{}
  >{\raggedright\arraybackslash}p{(\linewidth - 4\tabcolsep) * \real{0.2727}}
  >{\raggedright\arraybackslash}p{(\linewidth - 4\tabcolsep) * \real{0.4545}}
  >{\raggedright\arraybackslash}p{(\linewidth - 4\tabcolsep) * \real{0.2727}}@{}}
\toprule\noalign{}
\begin{minipage}[b]{\linewidth}\raggedright
目录
\end{minipage} & \begin{minipage}[b]{\linewidth}\raggedright
开放级别
\end{minipage} & \begin{minipage}[b]{\linewidth}\raggedright
说明
\end{minipage} \\
\midrule\noalign{}
\endhead
\bottomrule\noalign{}
\endlastfoot
\texttt{scx-research/core/} & 开源 & 数学框架核心代码 \\
\texttt{scx-research/adapters/} 中的 mlip/fem/vision/health 基础版 &
开源 & 领域适配器的基础实现 \\
\texttt{scx-research/adapters/semi/} 中的 toy demo & 开源 & 半导体 toy
demo 可开源 \\
\texttt{scx-research/adapters/music/} & 开源 & 探索性开源 \\
\texttt{scx-research/adapters/llm/} & 开源 & 仅框架层代码 \\
\texttt{scx-research/examples/} & 开源 & 所有示例 \\
\texttt{scx-pro/} & 闭源 & 商业级实现 \\
\texttt{experiments/} 中的原始大规模数据 & 视情况 & 部分受限 \\
\texttt{paper/} & 开源 & 论文全文 \\
\end{longtable}

\begin{center}\rule{0.5\linewidth}{0.5pt}\end{center}

\subsection{四、五篇论文谱系（2026-06-26
用户审定版）}\label{ux56dbux4e94ux7bc7ux8bbaux6587ux8c31ux7cfb2026-06-26-ux7528ux6237ux5ba1ux5b9aux7248}

\subsubsection{4.1 论文序列}\label{ux8bbaux6587ux5e8fux5217}

\begin{quote}
\textbf{评审意见}：GPT 原建议 SCX-MLIP → SCX-Theory → SCX-Sim →
SCX-Health。 \textbf{修正理由}：应该先 concrete 打根据地，再 abstract
建理论，再 apply back 验证理论，再跨领域扩展。 ``第一篇是 ACE
势函数合并（EGP 文件夹下），这是唯一有完整 DFT
训练数据的------根据地必须先拿下来。''
\end{quote}

\begin{longtable}[]{@{}
  >{\raggedright\arraybackslash}p{(\linewidth - 10\tabcolsep) * \real{0.1200}}
  >{\raggedright\arraybackslash}p{(\linewidth - 10\tabcolsep) * \real{0.1200}}
  >{\raggedright\arraybackslash}p{(\linewidth - 10\tabcolsep) * \real{0.1800}}
  >{\raggedright\arraybackslash}p{(\linewidth - 10\tabcolsep) * \real{0.2000}}
  >{\raggedright\arraybackslash}p{(\linewidth - 10\tabcolsep) * \real{0.2000}}
  >{\raggedright\arraybackslash}p{(\linewidth - 10\tabcolsep) * \real{0.1800}}@{}}
\toprule\noalign{}
\begin{minipage}[b]{\linewidth}\raggedright
顺序
\end{minipage} & \begin{minipage}[b]{\linewidth}\raggedright
论文
\end{minipage} & \begin{minipage}[b]{\linewidth}\raggedright
工作目录
\end{minipage} & \begin{minipage}[b]{\linewidth}\raggedright
核心问题
\end{minipage} & \begin{minipage}[b]{\linewidth}\raggedright
目标期刊
\end{minipage} & \begin{minipage}[b]{\linewidth}\raggedright
数据状态
\end{minipage} \\
\midrule\noalign{}
\endhead
\bottomrule\noalign{}
\endlastfoot
\textbf{1} & \textbf{EGP: Gauge-normalized residual expert merging for
ACE/PACE} & \texttt{egp/} & 多个 MLIP expert 在 ACE/PACE
下如何规范化、对齐、合并，避免 gauge ambiguity？ & npj Computational
Materials / PRM / JCTC & ✅ AlN v3 已训练，GaN/AlN v4 已提交超算 \\
\textbf{2} & \textbf{SCX-Theory} (State-Conditioned eXpertise) &
\texttt{SCX/} &
专家可靠性和数据价值应该被如何数学定义？什么时候全局专家排序失效？噪声
vs 可学习困难如何区分？ & TMLR / SIMODS / JMLR / arXiv 先占坑 & ✅
数学框架 + 合成实验（0 DFT） \\
\textbf{3} & \textbf{SCX-MLIP} (SCX applied to MLIP) & \texttt{SCX/} &
SCX 理论应用到 MLIP：state-conditioned expert reliability →
residual-state maps → error-guided DFT sampling → certified expert merge
& Nature Communications / npj Computational Materials & ⏳ 等 EGP Paper
1 + SCX-Theory 发表后 \\
\textbf{4} & \textbf{SCX-Sim} (Scientific \& Engineering Simulation) &
\texttt{SCX/} &
在高成本仿真中，如何根据状态选择粗模型/细模型/高保真求解器，降低成本并保持精度？
& Nature Computational Science / Nature Machine Intelligence & ⏳ toy
demo 可启动 \\
\textbf{5} & \textbf{SCX-Health} (Medical Data Valuation) &
\texttt{SCX/scx-health/} &
在医学数据中如何区分冗余样本、长尾高价值病例、疑似错标和专家依赖样本？ &
npj Digital Medicine & 🟡 MedMNIST 已下载，实验骨架已建 \\
\end{longtable}

\subsubsection{4.2 为什么 EGP Paper 1
必须排第一？}\label{ux4e3aux4ec0ux4e48-egp-paper-1-ux5fc5ux987bux6392ux7b2cux4e00}

\begin{enumerate}
\def\labelenumi{\arabic{enumi}.}
\tightlist
\item
  \textbf{唯一有完整 DFT 训练数据的}：AlN v3 425 frames 已训练，Model B
  gauge-fixed 已产出
\item
  \textbf{最
  concrete}：不是抽象数学，是有具体材料体系（AlN/GaN/AlGaN）的物理验证
\item
  \textbf{最容易被材料期刊审稿人理解}：gauge fixing 在 ACE
  线性框架下的意义清楚
\item
  \textbf{为 SCX-Theory
  提供''真实动机''}：审稿人会问''这个理论从哪里来？``→ 答：''来自 EGP
  Paper 1 的 gauge fixing 和 expert merging 实践经验''
\item
  \textbf{风险最低}：即使 SCX 理论后续被拒，至少有一篇材料方法学论文发表
\end{enumerate}

\subsubsection{4.3 为什么 SCX-Theory 排第二（而非 GPT
建议的排第一）？}\label{ux4e3aux4ec0ux4e48-scx-theory-ux6392ux7b2cux4e8cux800cux975e-gpt-ux5efaux8baeux7684ux6392ux7b2cux4e00}

GPT 建议 SCX-MLIP 排第一。但 SCX-MLIP 是''把 SCX 理论应用到
MLIP''------如果理论还没发表，应用论文就成了空中楼阁。

修正逻辑：

\begin{verbatim}
EGP Paper 1 (concrete, 材料方法)
    ↓ "我们在这个工作中发现 expert merging 需要 gauge fixing + state awareness"
SCX-Theory (abstract, 数学框架)  
    ↓ "我们把经验抽象为状态条件专家性理论"
SCX-MLIP (theory → application)
    ↓ "用 SCX 理论重新审视 MLIP expert merging，证明理论指导实践"
SCX-Sim + SCX-Health (跨领域扩展)
\end{verbatim}

EGP Paper 1 不需要引用 SCX-Theory（它是独立的 ACE 方法论文）。SCX-Theory
引用 EGP Paper 1 作为 motivation。SCX-MLIP 同时引用两者。

\subsubsection{4.4
每篇论文的独立贡献对象}\label{ux6bcfux7bc7ux8bbaux6587ux7684ux72ecux7acbux8d21ux732eux5bf9ux8c61}

\begin{longtable}[]{@{}
  >{\raggedright\arraybackslash}p{(\linewidth - 8\tabcolsep) * \real{0.1714}}
  >{\raggedright\arraybackslash}p{(\linewidth - 8\tabcolsep) * \real{0.1714}}
  >{\raggedright\arraybackslash}p{(\linewidth - 8\tabcolsep) * \real{0.2286}}
  >{\raggedright\arraybackslash}p{(\linewidth - 8\tabcolsep) * \real{0.2000}}
  >{\raggedright\arraybackslash}p{(\linewidth - 8\tabcolsep) * \real{0.2286}}@{}}
\toprule\noalign{}
\begin{minipage}[b]{\linewidth}\raggedright
论文
\end{minipage} & \begin{minipage}[b]{\linewidth}\raggedright
Loss
\end{minipage} & \begin{minipage}[b]{\linewidth}\raggedright
Expert
\end{minipage} & \begin{minipage}[b]{\linewidth}\raggedright
State
\end{minipage} & \begin{minipage}[b]{\linewidth}\raggedright
Action
\end{minipage} \\
\midrule\noalign{}
\endhead
\bottomrule\noalign{}
\endlastfoot
\textbf{EGP P1} & energy/force/stress residual & ACE experts
(shared+corrected) & atomic local environment (CN, bond, strain) & gauge
fix → merge \\
\textbf{SCX-Theory} & abstract loss ℓ & abstract experts f\_m &
measurable state partition S &
acquire/relabel/downweight/discard/route \\
\textbf{SCX-MLIP} & DFT energy/force residual & ACE/NEP/MACE & ACE
descriptor states & error-guided DFT sampling + certified merge \\
\textbf{SCX-Sim} & PDE/FEM residual & coarse/fine solver &
geometry/process state & high-fidelity solve / rerun \\
\textbf{SCX-Health} & classification error & model/doctor & visual
phenotype state & review/prune/retain \\
\end{longtable}

\subsubsection{4.5
发表顺序与策略红线}\label{ux53d1ux8868ux987aux5e8fux4e0eux7b56ux7565ux7ea2ux7ebf}

\begin{enumerate}
\def\labelenumi{\arabic{enumi}.}
\tightlist
\item
  \textbf{EGP Paper 1 优先}：已有 AlN v3 训练数据 + Model B，不需要等新
  DFT（等 GaN/AlN v4 数据回来后补充实验即可）。可在等超算期间写完
  Introduction + Methods + gauge fixing 理论 + AlN v3 实验结果
\item
  \textbf{SCX-Theory arXiv 占坑}：EGP Paper 1 投稿后立即发
  arXiv。不需要等任何 DFT------合成实验即可
\item
  \textbf{SCX-MLIP}：等 EGP Paper 1 接收 + SCX-Theory arXiv
  存在后，整合两者投稿
\item
  \textbf{SCX-Sim}：真正的 Nature 级目标，需 toy demo + 产业合作方数据
\item
  \textbf{SCX-Health}：开源影响力项目，MedMNIST 实验可立即跑，论文不急
\end{enumerate}

\textbf{红线：} - EGP Paper 1 不要塞 SCX 理论（只提 ``state-conditioned
expert reliability'' 作为 future work seed） - SCX-Theory 不要写太多
MLIP 结果（保持抽象，为 SCX-MLIP 留空间） - 不把 LLM 放主战场 -
每篇论文有独立贡献对象，不做 salami slicing

\begin{center}\rule{0.5\linewidth}{0.5pt}\end{center}

\subsection{五、LLM 专项}\label{ux4e94llm-ux4e13ux9879}

\subsubsection{5.1 GPT 对 LLM
的总体判断}\label{gpt-ux5bf9-llm-ux7684ux603bux4f53ux5224ux65ad}

\begin{quote}
\textbf{你的主战场不是大模型，而是科学计算与工程仿真。}
\end{quote}

核心逻辑： 1.
\textbf{大模型通用数据评价，大厂内部已经在做}，而且投入巨大 2.
\textbf{SCX 不具备与 GPT/Claude/DeepSeek 正面对打 judge 的资源} 3.
\textbf{LLM 数据缺少明确 y}*（无 gold label），SCX 的 claim 必须降级 4.
大厂强项是通用模型和工程规模；你的机会在垂直领域

\subsubsection{5.2 LLM 在 SCX
中的正确定位}\label{llm-ux5728-scx-ux4e2dux7684ux6b63ux786eux5b9aux4f4d}

LLM 在 SCX 中应作为：

\begin{quote}
\textbf{讨论 / extension / 远期外延，不是核心实验主战场。}
\end{quote}

SCX 在 LLM 中的正确切口不是''又一个 judge''，而是：

\begin{quote}
\textbf{状态条件评价器编排框架（State-conditioned judge orchestration
framework）}
\end{quote}

核心公式：

\begin{verbatim}
Q_SCX(x) = sum_j w_j(s(x)) * Q_j(x)
w_j(s) propto exp(alpha * A_j(s) - lambda * C_j)
\end{verbatim}

其中： - Q\_j(x)：第 j 个评价器的打分 - A\_j(s)：该评价器在状态 s
上的可靠性 - C\_j：评价成本

\subsubsection{5.3 SCX 与现有 9 类 LLM
数据筛选方法的区别}\label{scx-ux4e0eux73b0ux6709-9-ux7c7b-llm-ux6570ux636eux7b5bux9009ux65b9ux6cd5ux7684ux533aux522b}

\begin{longtable}[]{@{}
  >{\raggedright\arraybackslash}p{(\linewidth - 4\tabcolsep) * \real{0.2143}}
  >{\raggedright\arraybackslash}p{(\linewidth - 4\tabcolsep) * \real{0.3571}}
  >{\raggedright\arraybackslash}p{(\linewidth - 4\tabcolsep) * \real{0.4286}}@{}}
\toprule\noalign{}
\begin{minipage}[b]{\linewidth}\raggedright
方法
\end{minipage} & \begin{minipage}[b]{\linewidth}\raggedright
核心回答
\end{minipage} & \begin{minipage}[b]{\linewidth}\raggedright
SCX 的补充
\end{minipage} \\
\midrule\noalign{}
\endhead
\bottomrule\noalign{}
\endlastfoot
规则清洗 (FineWeb) & 是不是明显垃圾？ & 状态条件质量判断 \\
质量分类器 (FinerWeb) & 看起来是不是高质量？ &
多状态质量评估器，而非全局 \\
Perplexity/Loss (Rho-1) & 模型觉得难不难？ & state value + token value
结合 \\
目标分布匹配 (DSIR) & 像不像目标分布？ & 状态条件版 DSIR \\
Influence/LESS & 对目标能力有没有梯度贡献？ &
只在高风险/高价值状态调用 \\
LLM-as-a-Judge & 强模型觉得好不好？ & 多 judge 状态条件加权 \\
去重/多样性 & 是否重复？ & 质量 * (1-冗余) * 薄弱覆盖 \\
Verifier & 是否可验证正确？ & 设为高可靠专家，权重大 \\
DataComp-LM & 训练后效果如何？ & 作为最终验证闭环 \\
\end{longtable}

\subsubsection{5.4 SCX 的差异点}\label{scx-ux7684ux5deeux5f02ux70b9}

SCX 的核心差异不是''又一个质量打分算法''，而是：

\begin{quote}
\textbf{在当前状态 (s)、当前模型 (M)、当前数据集 (D)、可用评价器集合 (J)
下，这个样本应该采取什么动作？}
\end{quote}

动作空间：keep, drop, downweight, deduplicate, relabel, route to
stronger judge, verify with tool, use for SFT, use for preference data,
use for eval, use for curriculum later

\subsubsection{5.5 LLM 实验建议}\label{llm-ux5b9eux9a8cux5efaux8bae}

\textbf{不要一上来做大规模预训练。} 最小实验设计：

\begin{itemize}
\tightlist
\item
  \textbf{数据}：公开 instruction dataset（几千到几万条）
\item
  \textbf{模型}：0.5B--3B 小模型，LoRA 微调
\item
  \textbf{SCX 表示}：sentence embedding 或模型 hidden states
\item
  \textbf{误差代理}：初始模型 loss 或多模型 disagreement
\item
  \textbf{验证}：比较 full data vs random 20\% vs diversity 20\% vs
  high-loss 20\% vs SCX 20\%
\end{itemize}

\subsubsection{5.6 SCX-LLM
的商业定位}\label{scx-llm-ux7684ux5546ux4e1aux5b9aux4f4d}

\begin{itemize}
\tightlist
\item
  称为 \textbf{SCX Data Quality Calibration System}（数据质量校准系统）
\item
  不是''一个万能 judge''，而是''评价器的评价器''
\item
  类比：最牛 GPT 是医生，SCX 是医院质控体系
\item
  企业入口：数据打分 API、数据审计报告、私有 SCX-Judge 部署
\end{itemize}

\textbf{主要领域应聚焦垂直 judge（而非通用）：} -
SCX-MedJudge：医学问答/影像数据 - SCX-CodeJudge：代码数据 -
SCX-STEMJudge：数学/物理/材料科学

\begin{center}\rule{0.5\linewidth}{0.5pt}\end{center}

\subsection{六、附录：GPT
讨论中提到的关键结论}\label{ux516dux9644ux5f55gpt-ux8ba8ux8bbaux4e2dux63d0ux5230ux7684ux5173ux952eux7ed3ux8bba}

\subsubsection{6.1 SCX
的核心哲学}\label{scx-ux7684ux6838ux5fc3ux54f2ux5b66}

\begin{quote}
\textbf{评价质量决定数据质量，数据质量决定模型质量；但评价本身也是状态条件的，因此必须被校准、路由和审计。}
\end{quote}

\subsubsection{6.2 SCX
的适用条件（领域适配评估表）}\label{scx-ux7684ux9002ux7528ux6761ux4ef6ux9886ux57dfux9002ux914dux8bc4ux4f30ux8868}

\begin{longtable}[]{@{}lccc@{}}
\toprule\noalign{}
条件 & ACE/MLIP & 医学图像 & 大模型文本 \\
\midrule\noalign{}
\endhead
\bottomrule\noalign{}
\endlastfoot
有可靠标签/oracle & 强 & 中到强 & 弱到中 \\
表示空间有语义 & 强 & 中强 & 中 \\
状态可解释 & 强 & 中强 & 中弱 \\
噪声可验证 & 中强 & 中 & 弱 \\
冗余可识别 & 强 & 中强 & 中 \\
专家可靠性可估计 & 强 & 中强 & 中 \\
\end{longtable}

\subsubsection{6.3
噪声识别的理论边界（必须在论文中说明）}\label{ux566aux58f0ux8bc6ux522bux7684ux7406ux8bbaux8fb9ux754cux5fc5ux987bux5728ux8bbaux6587ux4e2dux8bf4ux660e}

\begin{quote}
\textbf{没有额外假设或参考信息时，高误差不能被唯一识别为噪声。}
\end{quote}

SCX 估计的是 noise risk（噪声风险），不是确定性噪声标签。论文中应包含
``Identifiability limits of noise detection'' 小节。

\subsubsection{6.4 SCX
的四种核心分类}\label{scx-ux7684ux56dbux79cdux6838ux5fc3ux5206ux7c7b}

\begin{longtable}[]{@{}
  >{\raggedright\arraybackslash}p{(\linewidth - 4\tabcolsep) * \real{0.2308}}
  >{\raggedright\arraybackslash}p{(\linewidth - 4\tabcolsep) * \real{0.3846}}
  >{\raggedright\arraybackslash}p{(\linewidth - 4\tabcolsep) * \real{0.3846}}@{}}
\toprule\noalign{}
\begin{minipage}[b]{\linewidth}\raggedright
类别
\end{minipage} & \begin{minipage}[b]{\linewidth}\raggedright
数学特征
\end{minipage} & \begin{minipage}[b]{\linewidth}\raggedright
建议动作
\end{minipage} \\
\midrule\noalign{}
\endhead
\bottomrule\noalign{}
\endlastfoot
Valuable（高价值） & r↑, ρ↑, C↑, A↑ & acquire / add data / train
expert \\
Noisy（疑似噪声） & r↑, ρ↓, C↓, A↓ & relabel / downweight / inspect \\
Redundant（冗余低价值） & r↓, ρ↑, Sim↑ & compress / weighted prototype /
prune \\
Expert-dependent（专家依赖） & R\_a(s) \textless\textless{} R\_b(s) &
route to best expert \\
\end{longtable}

\subsubsection{6.5 SCX 通用架构}\label{scx-ux901aux7528ux67b6ux6784}

\begin{verbatim}
x --[phi_domain]--> z --[Pi]--> s --> {value, redundancy, risk, expert, budget, model choice}
\end{verbatim}

\begin{itemize}
\tightlist
\item
  phi\_domain：领域专用状态编码器（Domain-Specific State Encoder /
  Gauge）
\item
  domain-invariant decision policy：通用数据动作策略
\end{itemize}

\subsubsection{6.6
最终战略定位}\label{ux6700ux7ec8ux6218ux7565ux5b9aux4f4d}

\begin{quote}
\textbf{SCX 不是替代求解器，不是替代医生，不是替代大模型；SCX
判断哪个专家、哪个仿真器、哪个数据样本在当前状态下可信和值得投入成本。}
\end{quote}

\begin{quote}
\textbf{SCX 不需要是万能方法；它是面向科学与工程仿真、医学
AI、高质量标注数据的状态条件质量判断与多保真调度框架。}
\end{quote}

\begin{center}\rule{0.5\linewidth}{0.5pt}\end{center}

\emph{本文件由 AI 助手根据 GPT 讨论记录（9011 行）整理，供 SCX
项目架构参考。}

\end{document}
